\documentclass[11pt,a4paper]{article}
\usepackage[T1]{fontenc}
\usepackage[utf8]{inputenc}
\usepackage[english]{babel}
\usepackage{lmodern}
\usepackage{microtype}
\usepackage{geometry}
\geometry{margin=2.5cm}
\usepackage{hyperref}
\hypersetup{
    colorlinks=true,
    linkcolor=blue!70!black,
    urlcolor=blue!70!black,
    citecolor=blue!70!black
}
\usepackage{booktabs}
\usepackage{longtable}
\usepackage{enumitem}
\usepackage{xcolor}
\usepackage{titlesec}
\usepackage{fancyhdr}
\usepackage{graphicx}
\usepackage{listings}

\definecolor{codegray}{rgb}{0.95,0.95,0.95}
\definecolor{doctum}{rgb}{0.10,0.30,0.60}

\lstset{
  backgroundcolor=\color{codegray},
  basicstyle=\small\ttfamily,
  breaklines=true,
  frame=single,
  rulecolor=\color{gray!40},
  xleftmargin=1em,
  framexleftmargin=0.5em
}

\pagestyle{fancy}
\fancyhf{}
\fancyhead[L]{\textcolor{doctum}{\small\textbf{Doctum Consilium}}}
\fancyhead[R]{\small\nouppercase{\leftmark}}
\fancyfoot[C]{\thepage}

\titleformat{\section}{\large\bfseries\color{doctum}}{}{0em}{}[\titlerule]
\titleformat{\subsection}{\normalsize\bfseries\color{doctum!80}}{}{0em}{}


\begin{document}

\title{
  \textcolor{doctum}{\Large\textbf{MedVision AI}} \\[0.5em]
  \large Medical Image AI Pipeline — Brain MRI Classification and Segmentation \\[0.3em]
  \normalsize Technical Foundations and Architecture
}
\author{Doctum Consilium \\ \small\textit{Generated 2026-05-08}}
\date{}
\maketitle
\thispagestyle{fancy}

\begin{abstract}
Deep learning pipeline for brain MRI analysis using transfer learning, DVC data versioning, MLflow experiment tracking, and Docker deployment.
This document covers the theoretical foundations, architectural decisions,
and key technical concepts required to understand, contribute to, and
maintain this repository.
\end{abstract}

\tableofcontents
\newpage

\section{Overview}

MedVision AI implements a reproducible medical image analysis pipeline that:
\begin{enumerate}
  \item Preprocesses and versions MRI datasets with DVC.
  \item Trains classification and segmentation models (DenseNet121, ResNet50, Swin-V2)
    using PyTorch with transfer learning from ImageNet weights.
  \item Tracks experiments, metrics, and artefacts with MLflow.
  \item Serves inference via a FastAPI REST endpoint.
\end{enumerate}
The pipeline targets brain tumour classification (glioma, meningioma, pituitary, no-tumour)
and is designed for extensibility to segmentation tasks (BraTS-compatible datasets).


\section{Architecture}

\begin{verbatim}
Raw MRI dataset (NIfTI/DICOM)
    │
    ▼
DVC pipeline (preprocess → train → evaluate)
    │                        │
    ▼                        ▼
Versioned artefacts     MLflow Tracking
    │                   (metrics, params, models)
    ▼
FastAPI inference endpoint
    │
    ▼
DICOM-SR / JSON prediction output
\end{verbatim}


\section{Key Technical Concepts}
\begin{itemize}[leftmargin=1.5em]
  \item Transfer Learning: fine-tuning ImageNet-pretrained CNNs on medical images; frozen backbone layers + trainable classification head; prevents overfitting on limited medical datasets.
  \item DenseNet121 architecture: dense connectivity $x_l = H_l([x_0, x_1, \ldots, x_{l-1}])$; feature reuse; proven on CheXNet (chest X-ray diagnosis).
  \item Data Augmentation: horizontal flip, rotation $\pm 10°$, colour jitter; critical for medical image generalisation given dataset scarcity.
  \item Class Imbalance: weighted sampling, macro-F1 as primary metric rather than accuracy to avoid misleading results on imbalanced classes.
  \item Reproducibility: DVC pins dataset versions + pipeline stages; seeds fixed across NumPy/PyTorch/CUDA for deterministic training.
  \item MLflow: automatic logging of hyperparameters, per-epoch metrics, confusion matrices, and model artifacts for experiment comparison.
\end{itemize}

\section{Data Flow}

DICOM/NIfTI → DVC preprocess stage (resize, normalise) →
PyTorch DataLoader (augmentation) →
Transfer learning model → loss + metrics →
MLflow log → checkpoint save →
FastAPI inference (\texttt{POST /predict}).


\section{Technology Stack}

\begin{tabular}{ll}
\toprule
\textbf{Component} & \textbf{Technology} \\
\midrule
Deep learning & PyTorch $\geq$ 2.0, torchvision \\
Model zoo & DenseNet121, ResNet50, Swin-V2-S \\
Data versioning & DVC 3.x \\
Experiment tracking & MLflow 2.x \\
Metrics & scikit-learn (F1, precision, recall, AUC) \\
Serving & FastAPI \\
\bottomrule
\end{tabular}


\section{Design Decisions}

\begin{itemize}
  \item DVC chosen over Git LFS for large dataset versioning with remote storage support.
  \item Macro-F1 as primary metric: clinically relevant — all classes matter equally.
  \item Three backbone options tested (DenseNet121, ResNet50, Swin-V2-S) to compare
    CNN vs.\ Vision Transformer performance on the same dataset.
\end{itemize}


\section{Repository Structure}
\begin{lstlisting}[language=bash,caption={Top-level directory layout}]
medvision-ai/
  .github/
    copilot-instructions.md   # GitHub Copilot guardrails
    agents/
      ecr-secrets-agent.agent.md  # ECR secret rotation runbook
  CLAUDE.md                   # Claude Code guardrails
  GEMINI.md                   # Gemini CLI guardrails
  INFRASTRUCTURE.md           # Operational runbook
  ROADMAP.md                  # Execution log and roadmap
  README.md                   # Project overview
\end{lstlisting}

\section{Contributing}
\begin{enumerate}
  \item Read \texttt{README.md}, \texttt{INFRASTRUCTURE.md}, and \texttt{ROADMAP.md} before any task.
  \item Follow the Code Documentation Standard: every public function must have
    a docstring (Google-style for Python, JSDoc for TypeScript, XML for C\#).
  \item Run the guardrails validator before committing:
    \begin{lstlisting}[language=bash]
git add -A && ./.guardrails/bin/validate_guardrails.sh
    \end{lstlisting}
  \item For deployment or destructive operations, always use a named tmux session:
    \begin{lstlisting}[language=bash]
tmux new-session -A -s deploy-medvision-ai
    \end{lstlisting}
\end{enumerate}

\end{document}
